% Reconstructed from the compiled lecture-note PDF.
\chapter{Linear algebra, matrix groups, and Lie theory}
This appendix develops the matrix language used in the proof. No abstract Lie theory is required beyond differentiating matrix products and exponentials.


\section{Two-dimensional determinants}

\begin{displaytext}
For column vectors v = (v1, v2)T and w = (w1, w2)T , write
\end{displaytext}

det(v, w) = v1w2 −v2w1.


Geometrically, det(v, w) is the signed area of the parallelogram spanned by v and w. It is positive when the ordered pair (v, w) is counterclockwise.


Let 0 −1! J = . 1 0


Then Jv is v rotated counterclockwise by π/2, and


det(v, w) = −vT Jw.


Indeed, 0 −1 ! w1 ! − v1 v2 = v1w2 −v2w1. 1 0 w2


If A is any 2 × 2 matrix, then


det(Av, Aw) = det(A) det(v, w).


This identity is the reason determinant-one matrices preserve oriented area.


\section{The group SL2(R)}

\begin{statementbox}{Definition}
Definition A.1. The special linear group is
\end{statementbox}

\begin{displaytext}
SL2(R) = \{g ∈M2(R) : det g = 1\} .
\end{displaytext}

It is a group under matrix multiplication.


A general element has the form


\begin{displaytext}
a b ! \
g = , ad −bc = 1, \
c d
\end{displaytext}

and d −b ! g−1 = . −c a


The equation ad−bc = 1 cuts out a three-dimensional surface in R4, so SL2(R) is a three-dimensional smooth manifold.


\subsection{Tangent vectors at the identity}

Let g(t) be a differentiable curve in SL2(R) with g(0) = I. Jacobi’s determinant formula gives


\begin{displaytext}
d \
0 = det g(t) = tr(g′(0)). \
dt t=0
\end{displaytext}

Thus a tangent vector at the identity must have trace zero.


\begin{statementbox}{Definition}
Definition A.2. The Lie algebra of SL2(R) is
\end{statementbox}

\begin{displaytext}
( a b ! ) \
sl2(R) = X = : a, b, c ∈R . \
c −a
\end{displaytext}

Its bracket is the commutator [X, Y ] = XY −Y X.


The word “Lie algebra” here simply means that sl2(R) is a vector space closed under commutators. The commutator records the first noncommutative correction in products of small matrix motions.


\section{The matrix exponential}

\begin{displaytext}
For any square matrix X, define \
∞ Xn \
exp(X) = X . \
n! n=0 \
The series converges absolutely. It satisfies
\end{displaytext}

\begin{displaytext}
d \
exp(tX) = X exp(tX) = exp(tX)X. \
dt
\end{displaytext}

\begin{displaytext}
If X and Y commute, then exp(X + Y ) = exp(X) exp(Y ). Without commutativity this formula \
generally fails. \
For X ∈sl2(R), the Cayley–Hamilton theorem gives an especially simple formula. Since tr X = 0, \
the characteristic polynomial is \
λ2 + det X.
\end{displaytext}

Therefore X2 = −det(X)I.


Write d = det X.


\begin{displaytext}
• If d > 0, then √ \
√ sin( d t) \
exp(tX) = cos( d t)I + √ X. \
d \
• If d = 0, then X2 = 0 and \
exp(tX) = I + tX. \
√ \
• If d < 0, put λ = −d. Then
\end{displaytext}

\begin{displaytext}
sinh(λt) \
exp(tX) = cosh(λt)I + X. \
λ
\end{displaytext}

These formulas make every constant-control segment in the Reinhardt system elementary.


\begin{statementbox}{Example}
Example A.3. Since J2 = −I,
\end{statementbox}

cos t −sin t ! exp(tJ) = . sin t cos t


Thus J is the infinitesimal generator of ordinary rotations.


\section{Left-trivialized velocity}

\begin{displaytext}
Let g(t) ∈SL2(R). Define \
X(t) = g(t)−1g′(t). \
Then X(t) ∈sl2(R), and the same equation can be written \
g′(t) = g(t)X(t).
\end{displaytext}

This is called the left-trivialized velocity because the tangent vector g′(t) is transported back to the identity by left multiplication with g(t)−1.


\begin{displaytext}
If v is a fixed vector and σ(t) = g(t)v, then
\end{displaytext}

\begin{displaytext}
σ′(t) = g′(t)v = g(t)X(t)v.
\end{displaytext}

Thus one matrix X(t) describes the simultaneous velocities of all vectors moved by g(t).


If g is reparameterized by t = t(s), then


\begin{displaytext}
dt \
eX(s) = X(t(s)) ds. \
Consequently det eX = (dt/ds)2 det X. When det X > 0, one may choose the parameter so that \
det X ≡1.
\end{displaytext}

\section{Adjoint action and Lax equations}

\begin{displaytext}
For g ∈SL2(R) and X ∈sl2(R), define \
Adg X = gXg−1.
\end{displaytext}

This is conjugation. It preserves trace, determinant, eigenvalues, and commutators:


Adg[X, Y ] = [Adg X, Adg Y ].


\begin{displaytext}
Differentiating conjugation gives the basic Lax identity. Let P(t) and X(t) be matrix curves. If \
X(t) = h(t)X0h(t)−1, h′(t) = P(t)h(t),
\end{displaytext}

\begin{displaytext}
then \
X′ = [P, X]. \
Conversely, a solution of X′ = [P, X] evolves by conjugation. Therefore its characteristic polynomial \
is constant. \
Proposition A.4 (Isospectrality). If X′ = [P, X], then tr(Xk) is constant for every positive integer \
k. In particular, for 2 × 2 trace-zero matrices, det X is constant.
\end{displaytext}

\begin{prooftext}
Proof. By cyclicity of trace,
\end{prooftext}

\begin{displaytext}
d \
tr(Xk) = k tr(Xk−1[P, X]) = k tr(Xk−1PX −XkP) = 0. \
dt
\end{displaytext}

The Reinhardt state equation is a control-dependent Lax equation: Zu X′ = ⟨Zu, X⟩, X .


\section{The trace pairing}

On sl2(R), define ⟨X, Y ⟩= tr(XY ).


This is a symmetric bilinear form. It is nondegenerate but not positive definite. For a b ! X = , c −a


\begin{displaytext}
one has \
⟨X, X⟩= 2(a2 + bc) = −2 det X.
\end{displaytext}

Thus matrices with positive determinant have negative squared length in this indefinite geometry.


\begin{displaytext}
The pairing is invariant under conjugation: \
⟨Adg X, Adg Y ⟩= tr(gXY g−1) = tr(XY ).
\end{displaytext}

It also satisfies ⟨X, [Y, Z]⟩= ⟨[X, Y ], Z⟩.


This follows from expanding and cyclically permuting traces.


\subsection{Useful 2 × 2 identities}

\begin{displaytext}
For X, Y ∈sl2(R), \
XY + Y X = ⟨X, Y ⟩I. \
To verify it, write out the entries or polarize X2 = −det(X)I. Several consequences are used \
repeatedly: \
[X, Y ]2 = ⟨X, Y ⟩2 −⟨X, X⟩⟨Y, Y ⟩ I, \
⟨[X, Y ], [X, Y ]⟩= 2 ⟨X, Y ⟩2 −⟨X, X⟩⟨Y, Y ⟩ , \
[X, [X, Y ]] = 2 ⟨X, X⟩Y −2 ⟨X, Y ⟩X. \
These are the indefinite analogues of vector triple-product identities in R3.
\end{displaytext}

\section{A concrete basis for sl2(R)}

\begin{displaytext}
It is useful to use \
1 0 ! 0 1 ! 0 0 ! \
H = , E = , F = . \
0 −1 0 0 1 0
\end{displaytext}

\begin{displaytext}
Then \
[H, E] = 2E, [H, F] = −2F, [E, F] = H. \
Every X ∈sl2(R) is aH + bE + cF.
\end{displaytext}

Another basis adapted to geometry is


\begin{displaytext}
0 −1! 0 1 ! 1 0 ! \
J = , K = , H = . \
1 0 1 0 0 −1
\end{displaytext}

The trace pairing has signs


\begin{displaytext}
⟨J, J⟩= −2, ⟨K, K⟩= 2, ⟨H, H⟩= 2,
\end{displaytext}

and the basis vectors are pairwise orthogonal. Hence sl2(R) is a copy of three-dimensional Minkowski space with signature (−, +, +).


\section{Adjoint orbits as hyperboloids}

Because det J = 1, the adjoint orbit


\begin{displaytext}
n o OJ = gJg−1 : g ∈SL2(R)
\end{displaytext}

\begin{displaytext}
lies in the quadratic surface det X = 1. Since ⟨X, X⟩= −2 det X, this is one sheet of a two-sheeted \
hyperboloid in the Minkowski coordinates above. The sign condition ⟨J, X⟩< 0 chooses the sheet \
containing J. \
The stabilizer of J is the rotation group SO2(R). Therefore \
OJ ∼= SL2(R)/SO2(R).
\end{displaytext}

The same quotient is the hyperbolic plane. This is the structural reason hyperbolic geometry appears in the control problem.


\section{The upper-half-plane parametrization}

For z = x + iy with y > 0, define


x/y −(x2 + y2)/y ! Φ(z) = . 1/y −x/y


A direct calculation gives


\begin{displaytext}
+ y2 + 1 tr Φ(z) = 0, det Φ(z) = 1, ⟨J, Φ(z)⟩= −x2 < 0. \
y
\end{displaytext}

Thus Φ(z) ∈OJ. The group SL2(R) acts on the upper half-plane by Mobius transformations:


a b ! az + b · z = c d cz + d.


The map Φ is equivariant: Φ(g · z) = gΦ(z)g−1.


One can verify this directly, though the quotient description makes it conceptually inevitable.


\section{Differentiating determinants: Jacobi’s formula}

\begin{statementbox}{Theorem}
Theorem A.5 (Jacobi). For a differentiable invertible matrix curve A(t),
\end{statementbox}

d det A(t) = det A(t) tr A(t)−1A′(t) . dt


For a 2 × 2 matrix, the formula can be checked from coordinates. It implies:


\begin{displaytext}
• if g(t) ∈SL2(R), then g−1g′ has trace zero; \
• if X(t) has constant determinant, then tr(X−1X′) = 0; \
• the normalization det X = 1 is compatible with the state equation.
\end{displaytext}

\section{The Cayley transform and su(1, 1)}

\begin{displaytext}
The source sometimes changes from real matrices to complex matrices because rotations become \
multiplication by a phase. Set \
1 1 −i! \
C = √ . \
2 1 i
\end{displaytext}

Conjugation X 7→CXC−1 identifies sl2(R) with


\begin{displaytext}
( it z ! ) \
su(1, 1) = : t ∈R, z ∈C . \
¯z −it
\end{displaytext}

For a b ! X = , c −a


the transformed coordinates can be written


\begin{displaytext}
it z ! b + c b −c \
CXC−1 = , z = + ia, t = . \
¯z −it 2 2
\end{displaytext}

The determinant condition becomes


t2 −|z|2 = det X.


For determinant 1, this is the upper sheet of a hyperboloid when t > 0. These are the hyperboloid coordinates used near the singular locus.


\section{Matrix product rules used in the proof}

The following elementary formulas are worth memorizing:


\begin{displaytext}
(AB)′ = A′B + AB′, \
(A−1)′ = −A−1A′A−1, \
(ABA−1)′ = A′BA−1 + AB′A−1 −ABA−1A′A−1, \
d \
⟨A, B⟩= A′, B + A, B′ . \
dt
\end{displaytext}

If A′ = AP, the conjugation formula simplifies to


(ABA−1)′ = A(B′ + [P, B])A−1.


\section{Exercises}

\begin{statementbox}{Exercise}
Exercise A.6. Prove XY + Y X = tr(XY )I for all X, Y ∈sl2(R) by direct multiplication.
\end{statementbox}

\begin{statementbox}{Exercise}
Exercise A.7. Show that X′ = [P, X] implies det X is constant without using the general tracepower argument.
\end{statementbox}

\begin{statementbox}{Exercise}
Exercise A.8. Verify explicitly that Φ(i) = J and that Φ(g · i) = gJg−1.
\end{statementbox}

\begin{statementbox}{Exercise}
Exercise A.9. Let Z be constant and suppose X′ = [Z, X]/ ⟨Z, X⟩. Prove that ⟨Z, X⟩is constant and obtain Z X(t) = exp(tP)X(0) exp(−tP), P = ⟨Z, X(0)⟩.
\end{statementbox}

\begin{insightbox}{Chapter summary}
\end{insightbox}

All Lie-theoretic operations in the proof reduce to explicit 2×2 matrix calculus. The essential structures are determinant-one matrices, trace-zero velocities, conjugation, the invariant trace


pairing, and the identification of the determinant-one adjoint orbit with the hyperbolic plane.

